
%\needspace{4\baselineskip}
\color{uniblue}{\subsection{Полуавтоматическое включение (ПАВ)}}
\color{black}

\begin{enumerate}[label=\arabic{section}.\arabic{subsection}.\arabic{enumi}, labelsep=4pt, leftmargin=0pt, itemindent=57pt, itemsep=0pt, parsep=0pt] % если что можно крутить itemsep

\item
В устройстве предусмотрен режим полуавтоматического включения (ПАВ). Данная функция предполагает включение выключателя с контролем синхронизма или улавливанием синхронизма при условии оперативного включения выключателя на противоположном конце линии. Для ввода режима ПАВ В1 предусмотрен входной сигнал ФСУ \textit{«ПАВ В1»}. При оперативном включении линии с противоположного конца и выполнении условий синхронизма происходит включение выключателя. Если в течение времени \textbf{«Тожид ПАВ В1»} не будут выполнены условия синхронизма, то схема ПАВ В1 выведется из работы.

\item
Функциональная схема логики ПАВ В1 приведена на рисунке \ref{pict_pav:PAV}. Функциональная схема логики ПАВ В2 выполняется аналогично схеме логики ПАВ В1.

\vspace{2mm}
\begin{figure}[H]
\centering
\includegraphics[width=0.9\textwidth,height=0.9\textheight,keepaspectratio]{\fbpath/761.3061 ПАВ/img/PAV.pdf}
\captionof{figure}{Функциональная схема логики ПАВ В1}\label{pict_pav:PAV}
\end{figure}

\item
Уставки ПАВ В1 представлены в таблице \ref{desc_apv:tbl1}. Уставки ПАВ В2 аналогичны уставкам ПАВ В1.

\small
\begin{longtable}{|>{\centering\arraybackslash}m{5.3cm}|>{\centering\arraybackslash}m{3.3cm}|>{\centering\arraybackslash}m{4.2cm}|>{\centering\arraybackslash}m{1.8cm}|>{\centering\arraybackslash}m{1cm}|}
\caption{Параметры для настройки ПАВ В1\hfill\vspace{-0.5\baselineskip}}\label{desc_apv:tbl1}\\ 
\hline
\rowcolor{gray!30}
Параметр & Условное обозначение на схеме & Значение/ Диапазон & Единица измерения & Шаг \\ 
\hline
\endfirsthead
\caption*{\hspace{3pt}\emph{Продолжение таблицы \ref{desc_apv:tbl1}\hfill\vspace{-0.5\baselineskip}}} \\ % сделано по ГОСТ 2.105 п.6.8.7
\hline
\rowcolor{gray!30}
Параметр & Условное обозначение на схеме & Значение/ Диапазон & Единица измерения & Шаг \\ 
%\hline
\endhead

%===m>ПАВ В1|ПАВ
\raggedright \textcolor{red}{FullDescription НЕ НАЙДЕН!} & \centering ПАВ & \centering Введено \\ Выведено & \centering --- & \centering \arraybackslash --- \\
\hline
\raggedright \textcolor{red}{FullDescription НЕ НАЙДЕН!} & \centering Тожид ПАВ & \centering 0 ... 1800000 & \centering мс & \centering \arraybackslash 5 \\
\hline
\raggedright \textcolor{red}{FullDescription НЕ НАЙДЕН!} & \centering Тимп ПАВ & \centering 0 ... 300000 & \centering мс & \centering \arraybackslash 5 \\
\hline
%===m
\end{longtable}
\normalsize

\end{enumerate}

\newpage